% Converted from the original powerdot Jeopardy slides (Calculus II).
% The worked solutions from the slides are in \solution; the \response is
% the final answer.
\documentclass{article}
% Which PDF to make (the web game ignores this):
%   questions  clues only, e.g. a study sheet for students
%   answers    clues and responses
%   solutions  clues, responses, and worked solutions
\usepackage[solutions]{../jeopardy}

\gametitle{Calculus II Jeopardy}

\begin{document}

%======================================================================
\begin{category}{Integration by Parts}

  \begin{clue}
    \[ \int x e^{x} \; dx \]
  \response
    What is $x e^{x} - e^{x} + c$?
  \solution
    Use integration by parts with $u=x$ and $dv=e^{x}\,dx.$ Then $du=dx$ and $v=e^{x}.$ So,
    \begin{align*}
    \int x e^{x} \; dx
    &= x e^{x} - \int e^{x} \; dx \\
    &= x e^{x} - e^{x} + c.
    \end{align*}
  \end{clue}

  \begin{clue}
    \[ \int x \ln x \; dx \]
  \response
    What is $\dfrac{x^{2}}{2} \ln x - \dfrac{1}{4} x^{2} + c$?
  \solution
    Use integration by parts with $u=\ln x$ and $dv=x\,dx.$ Then $du=\frac{1}{x}\,dx$ and $v=\frac{x^{2}}{2}.$ So,
    \begin{align*}
    \int x \ln x \; dx
    &= \frac{x^{2}}{2} \ln x - \int \frac{x^{2}}{2x} \; dx \\
    &= \frac{x^{2}}{2} \ln x - \frac{1}{2} \int x \; dx \\
    &= \frac{x^{2}}{2} \ln x - \frac{1}{4} x^{2} + c.
    \end{align*}
  \end{clue}

  \begin{clue}
    \[ \int (x+2) e^{2x+1} \; dx \]
  \response
    What is $\dfrac{1}{2}e^{2x+1}(x+2) - \dfrac{1}{4}e^{2x+1} + c$?
  \solution
    Use integration by parts with $u=x+2$ and $dv=e^{2x+1}\,dx.$ Then $du=dx$ and $v=\frac{1}{2}e^{2x+1}.$ So,
    \begin{align*}
    \int (x+2) e^{2x+1} \; dx
    &= \frac{1}{2}e^{2x+1}(x+2) - \int \frac{1}{2}e^{2x+1} \; dx \\
    &= \frac{1}{2}e^{2x+1}(x+2) - \frac{1}{4}e^{2x+1} + c.
    \end{align*}
  \end{clue}

  \begin{clue}
    \[ \int \sin^{-1} x \; dx \]
  \response
    What is $x\sin^{-1} x + \sqrt{1-x^{2}} + c$?
  \solution
    Use integration by parts with $u=\sin^{-1} x$ and $dv=dx.$ Then $du=\frac{1}{\sqrt{1-x^{2}}}\,dx$ and $v=x.$ So,
    \[
    \int \sin^{-1} x \; dx = x\sin^{-1} x - \int x (1-x^{2})^{-\frac{1}{2}} \; dx.
    \]
    For the last integral, substitute $w=1-x^{2},$ so $dw=-2x\,dx.$ Then
    \begin{align*}
    \int \sin^{-1} x \; dx
    &= x\sin^{-1} x + \frac{1}{2} \int w^{-\frac{1}{2}} \; dw \\
    &= x\sin^{-1} x + w^{\frac{1}{2}} + c \\
    &= x\sin^{-1} x + \sqrt{1-x^{2}} + c.
    \end{align*}
  \end{clue}

  \begin{clue}
    \[ \int x^{2} \sin x \; dx \]
  \response
    What is $-x^{2} \cos x + 2x \sin x + 2 \cos x + c$?
  \solution
    Use integration by parts with $u=x^{2}$ and $dv=\sin x\,dx.$ Then $du=2x\,dx$ and $v=-\cos x.$ So,
    \[
    \int x^{2} \sin x \; dx = -x^{2} \cos x + \int 2x \cos x \; dx.
    \]
    Use integration by parts again, with $u=2x$ and $dv=\cos x\,dx.$ Then $du=2\,dx$ and $v=\sin x.$ So,
    \begin{align*}
    \int x^{2} \sin x \; dx
    &= -x^{2} \cos x + \left( 2x \sin x - \int 2 \sin x \; dx \right) \\
    &= -x^{2} \cos x + 2x \sin x + 2 \cos x + c.
    \end{align*}
  \end{clue}

\end{category}

%======================================================================
\begin{category}{Limits}

  \begin{clue}
    This is the value $\displaystyle\lim_{x\to\infty} e^{-x^2}.$
  \response
    What is $0$?
  \end{clue}

  \begin{clue}
    This is the value $\displaystyle \lim_{x\to 0}\frac{x-\sin x}{x^2}.$
  \response
    What is $0$?
  \solution
    \[
    \begin{array}{rcl}
    {\displaystyle \lim_{x\to 0}\frac{x-\sin x}{x^2}} & = & {\displaystyle \lim_{x\to 0}\frac{1-\cos x}{2x}} \\
    & = & {\displaystyle \lim_{x\to 0}\frac{\sin x}{2}}\\
    & = & 0
    \end{array}
    \]
  \end{clue}

  \begin{clue}
    This is the value $\displaystyle\lim_{x\to\infty}\frac{\log_5 x}{\log_6 x}.$
  \response
    What is $\dfrac{\ln 6}{\ln 5}$?
  \solution
    \[
    \begin{array}{rcl}
    {\displaystyle\lim_{x\to\infty}\frac{\log_5 x}{\log_6 x}} & = & {\displaystyle\lim_{x\to\infty}\frac{\Big(\frac{1}{x\ln 5}\Big)}{\Big(\frac{1}{x\ln 6}\Big)}} \\
    & = & {\displaystyle\lim_{x\to\infty}\frac{\ln 6}{\ln 5}}\\
    & = & {\displaystyle \frac{\ln 6}{\ln 5}}
    \end{array}
    \]
  \end{clue}

  \begin{clue}
    This is the value $\displaystyle \lim_{x\to 0}(3\csc x-3\cot x).$
  \response
    What is $0$?
  \solution
    \[
    \begin{array}{rcl}
    {\displaystyle \lim_{x\to 0}(3\csc x-3\cot x)} & = & {\displaystyle \lim_{x\to 0}\Big(\frac{3}{\sin x}-\frac{3\cos x}{\sin x}\Big)} \\
    & = & {\displaystyle \lim_{x\to 0}\Big(\frac{3-3\cos x}{\sin x}\Big)}\\
    & = & {\displaystyle \lim_{x\to 0}\Big(\frac{3\sin x}{\cos x}\Big)}\\
    & = & 0
    \end{array}
    \]
  \end{clue}

  \begin{clue}
    This is the value $\displaystyle\lim_{x\to 0^{+}}x^x.$
  \response
    What is $1$?
  \solution
    \[
    \begin{array}{rcl}
    {\displaystyle \lim_{x\to 0^{+}}\ln (x^{x})} & = & {\displaystyle \lim_{x\to 0^{+}}x\ln x} \\
    & = & {\displaystyle \lim_{x\to 0^{+}}\frac{\ln x}{x^{-1}}}\\
    & = & {\displaystyle \lim_{x\to 0^{+}}\frac{x^{-1}}{-x^{-2}}}\\
    & = & {\displaystyle \lim_{x\to 0^{+}}-x}\\
    & = & 0\\
    {\displaystyle \lim_{x\to 0^{+}} x^{x}} & = & e^{0}=1
    \end{array}
    \]
  \end{clue}

\end{category}

%======================================================================
\begin{category}{Proper or Improper?}

  \begin{clue}
    \[ \int_{-3}^{\infty} \frac{x^{3}}{7+ \sin^{2} x} \, dx \]
  \response
    What is IMPROPER?
  \end{clue}

  \begin{clue}
    \[ \int_{-3}^{10} \frac{x^{3}}{(x+5)(x-11)} \, dx \]
  \response
    What is PROPER?
  \end{clue}

  \begin{clue}
    \[ \int_{-\pi/2}^{\pi/2} \frac{e^{x}}{\cos x} \, dx \]
  \response
    What is IMPROPER?
  \end{clue}

  \begin{clue}
    \[ \int_{\pi/6}^{\pi/4} \frac{1}{\sec x} \, dx \]
  \response
    What is PROPER?
  \end{clue}

  \begin{clue}
    Wearing white shoes between Labor Day and Easter
  \response
    What is IMPROPER?
  \end{clue}

\end{category}

%======================================================================
\begin{category}{Applied Integration}

  \begin{clue}
    This is the INTEGRAL used to compute the area bounded by the curves
    $y=x^2+1,$ $y=2x$ and $x=0.$
  \response
    What is $\displaystyle A=\int_0^1 (x^2+1-2x)\; dx$?
  \end{clue}

  \begin{clue}
    This is the INTEGRAL used to compute the volume of the solid obtained by
    revolving the region bounded by the curves $y=x^2+1,$ $y=2x$ and $x=0$
    about the $x$-axis.
  \response
    What is
    \[V=\pi\int_0^1 \left((x^2+1)^2-(2x)^2\right) dx\,?\]
  \end{clue}

  \begin{clue}[dailydouble]
    This is the INTEGRAL used to compute the volume of the solid obtained by
    revolving the region bounded by the curves $y=x^2+1,$ $y=2x$ and $x=0$
    about the line $x=-2.$
  \response
    What is
    \[V=2\pi\int_0^1 (x+2)(x^2+1-2x)\, dx\,?\]
  \end{clue}

  \begin{clue}
    These are the INTEGRALS needed to find the centroid of the region bounded
    by $y=\sin x$ and the $x$-axis on $0\le x\le \pi.$
  \response
    What are
    \[\bar x=\frac{1}{A}\int_{0}^{\pi}x\sin x\; dx\]
    and
    \[\bar y=\frac{1}{2A}\int_{0}^{\pi}\sin^{2}x\;dx,\]
    where $\displaystyle A=\int_{0}^{\pi}\sin x\; dx$?
  \end{clue}

  \begin{clue}
    This is the INTEGRAL used to compute the arclength of one period of the
    graph of $y=\sin x.$
  \response
    What is $\displaystyle L=\int_0^{2\pi}\sqrt{1+\cos^2 x}\; dx$?
  \end{clue}

\end{category}

%======================================================================
\begin{category}{Creative Calculus}

  \begin{clue}
    \[ \int \left[ \frac{d}{dx} f(x) \right] \; dx \]
  \response
    What is $f(x) + c$?
  \end{clue}

  \begin{clue}
    \[ \int \cos t \; dx \]
  \response
    What is $(\cos t)x + c$?
  \solution
    \[
    \begin{array}{rcl}
    {\displaystyle \int \cos t \; dx} & = & {\displaystyle \cos t \int  dx} \\
     & = & {\displaystyle (\cos t)x +   c }
    \end{array}
    \]
  \end{clue}

  \begin{clue}
    \[ \frac{d}{dr} \int_{3}^{r} \sin (x^{2}) \; dx \]
  \response
    What is $\sin (r^{2})$?
  \end{clue}

  \begin{clue}
    \[ \int \frac{\cos x}{1 + \sin^{2} x} \; dx \]
  \response
    What is $\tan^{-1} (\sin x) + c$?
  \solution
    \[
    \begin{array}{rcl}
    {\displaystyle \int \frac{\cos x}{1 + \sin^{2} x} \; dx} & = & {\displaystyle \int \frac{du}{1+u^{2}}} \\
    \left( \begin{array}{rcl}
    u & = & \sin x \\
    du & = & \cos x \; dx
    \end{array} \right)^{\nearrow} & = & {\displaystyle \tan^{-1} u + c} \\
     & = & {\displaystyle \tan^{-1} (\sin x) + c}
    \end{array}
    \]
  \end{clue}

  \begin{clue}
    \[ \int e^{x+e^{x}} \; dx \]
  \response
    What is $e^{e^{x}} + c$?
  \solution
    \[
    \begin{array}{rcl}
    {\displaystyle \int e^{x+e^{x}} \; dx} & = & {\displaystyle \int e^{x}e^{e^{x}} \; dx} \\
    \left( \begin{array}{rcl}
    u & = & e^{x} \\
    du & = & e^{x} \; dx
    \end{array} \right) \rightarrow & = & {\displaystyle \int e^{u} \; du} \\
     & = & {\displaystyle e^{u} + c} \\
     & = & {\displaystyle e^{e^{x}} + c }
    \end{array}
    \]
  \end{clue}

\end{category}

%======================================================================
\begin{category}{Potpourri}

  \begin{clue}
    An important player in the creation of calculus other than Newton and Leibniz.
  \response
    Who is Fermat? (Many questions are possible, e.g., Barrow or Cavalieri.)
  \end{clue}

  \begin{clue}
    This is the name of the new science complex.
  \response
    What is the Rector Science Complex?
  \end{clue}

  \begin{clue}
    The smallest of the irrational numbers $\pi,$ $e,$ and $\sqrt{2}.$
  \response
    What is $\sqrt{2}$?
  \end{clue}

  \begin{clue}
    The second full moon in a single month
  \response
    What is a blue moon?
  \end{clue}

  \begin{clue}
    A ``mathematical'' car from Acura

    (at least the ``LS'' model)
  \response
    What is the ``INTEGRA LS''?
  \end{clue}

\end{category}

%======================================================================
\begin{final}{Musical instruments}
  This is the volume of \emph{Gabriel's horn}, i.e., the solid obtained by
  revolving $\displaystyle y=\frac{1}{x}$ (for $1\le x<\infty$) about the
  $x$-axis.
\response
  What is $\pi$?
\solution
  \[
  \begin{array}{rcl}
  {\displaystyle \int_{1}^{\infty}\pi\Big(\frac{1}{x}\Big)^{2}\; dx} & = & {\displaystyle \lim_{b\to \infty} \int_{1}^{b}\frac{\pi}{x^{2}}\; dx} \\
  & = & {\displaystyle \lim_{b\to \infty}\Big(-\frac{\pi}{x}\Big)\Big]_{1}^{b}}\\
  & = & {\displaystyle \lim_{b\to \infty}\Big(-\frac{\pi}{b}+\frac{\pi}{1}\Big)}\\
  & = & \pi
  \end{array}
  \]
\end{final}

\end{document}
